\IfFileExists{stacks-project.cls}{%
\documentclass{stacks-project}
}{%
\documentclass{amsart}
}

% For dealing with references we use the comment environment
\usepackage{verbatim}
\newenvironment{reference}{\comment}{\endcomment}
%\newenvironment{reference}{}{}
\newenvironment{slogan}{\comment}{\endcomment}
\newenvironment{history}{\comment}{\endcomment}

% For commutative diagrams we use Xy-pic
\usepackage[all]{xy}

% We use 2cell for 2-commutative diagrams.
\xyoption{2cell}
\UseAllTwocells

% We use multicol for the list of chapters between chapters
\usepackage{multicol}

% This is generally recommended for better output
\usepackage{lmodern}
\usepackage[T1]{fontenc}

% For cross-file-references
\usepackage{xr-hyper}

% Package for hypertext links:
\usepackage{hyperref}

% For any local file, say "hello.tex" you want to link to please
% use \externaldocument[hello-]{hello}
\externaldocument[introduction-]{introduction}
\externaldocument[conventions-]{conventions}
\externaldocument[sets-]{sets}
\externaldocument[categories-]{categories}
\externaldocument[topology-]{topology}
\externaldocument[sheaves-]{sheaves}
\externaldocument[sites-]{sites}
\externaldocument[stacks-]{stacks}
\externaldocument[fields-]{fields}
\externaldocument[algebra-]{algebra}
\externaldocument[brauer-]{brauer}
\externaldocument[homology-]{homology}
\externaldocument[derived-]{derived}
\externaldocument[simplicial-]{simplicial}
\externaldocument[more-algebra-]{more-algebra}
\externaldocument[smoothing-]{smoothing}
\externaldocument[modules-]{modules}
\externaldocument[sites-modules-]{sites-modules}
\externaldocument[injectives-]{injectives}
\externaldocument[cohomology-]{cohomology}
\externaldocument[sites-cohomology-]{sites-cohomology}
\externaldocument[dga-]{dga}
\externaldocument[dpa-]{dpa}
\externaldocument[sdga-]{sdga}
\externaldocument[hypercovering-]{hypercovering}
\externaldocument[schemes-]{schemes}
\externaldocument[constructions-]{constructions}
\externaldocument[properties-]{properties}
\externaldocument[morphisms-]{morphisms}
\externaldocument[coherent-]{coherent}
\externaldocument[divisors-]{divisors}
\externaldocument[limits-]{limits}
\externaldocument[varieties-]{varieties}
\externaldocument[topologies-]{topologies}
\externaldocument[descent-]{descent}
\externaldocument[perfect-]{perfect}
\externaldocument[more-morphisms-]{more-morphisms}
\externaldocument[flat-]{flat}
\externaldocument[groupoids-]{groupoids}
\externaldocument[more-groupoids-]{more-groupoids}
\externaldocument[etale-]{etale}
\externaldocument[chow-]{chow}
\externaldocument[intersection-]{intersection}
\externaldocument[pic-]{pic}
\externaldocument[weil-]{weil}
\externaldocument[adequate-]{adequate}
\externaldocument[dualizing-]{dualizing}
\externaldocument[duality-]{duality}
\externaldocument[discriminant-]{discriminant}
\externaldocument[derham-]{derham}
\externaldocument[local-cohomology-]{local-cohomology}
\externaldocument[algebraization-]{algebraization}
\externaldocument[curves-]{curves}
\externaldocument[resolve-]{resolve}
\externaldocument[models-]{models}
\externaldocument[functors-]{functors}
\externaldocument[equiv-]{equiv}
\externaldocument[pione-]{pione}
\externaldocument[etale-cohomology-]{etale-cohomology}
\externaldocument[proetale-]{proetale}
\externaldocument[relative-cycles-]{relative-cycles}
\externaldocument[more-etale-]{more-etale}
\externaldocument[trace-]{trace}
\externaldocument[crystalline-]{crystalline}
\externaldocument[spaces-]{spaces}
\externaldocument[spaces-properties-]{spaces-properties}
\externaldocument[spaces-morphisms-]{spaces-morphisms}
\externaldocument[decent-spaces-]{decent-spaces}
\externaldocument[spaces-cohomology-]{spaces-cohomology}
\externaldocument[spaces-limits-]{spaces-limits}
\externaldocument[spaces-divisors-]{spaces-divisors}
\externaldocument[spaces-over-fields-]{spaces-over-fields}
\externaldocument[spaces-topologies-]{spaces-topologies}
\externaldocument[spaces-descent-]{spaces-descent}
\externaldocument[spaces-perfect-]{spaces-perfect}
\externaldocument[spaces-more-morphisms-]{spaces-more-morphisms}
\externaldocument[spaces-flat-]{spaces-flat}
\externaldocument[spaces-groupoids-]{spaces-groupoids}
\externaldocument[spaces-more-groupoids-]{spaces-more-groupoids}
\externaldocument[bootstrap-]{bootstrap}
\externaldocument[spaces-pushouts-]{spaces-pushouts}
\externaldocument[spaces-chow-]{spaces-chow}
\externaldocument[groupoids-quotients-]{groupoids-quotients}
\externaldocument[spaces-more-cohomology-]{spaces-more-cohomology}
\externaldocument[spaces-simplicial-]{spaces-simplicial}
\externaldocument[spaces-duality-]{spaces-duality}
\externaldocument[formal-spaces-]{formal-spaces}
\externaldocument[restricted-]{restricted}
\externaldocument[spaces-resolve-]{spaces-resolve}
\externaldocument[formal-defos-]{formal-defos}
\externaldocument[defos-]{defos}
\externaldocument[cotangent-]{cotangent}
\externaldocument[examples-defos-]{examples-defos}
\externaldocument[algebraic-]{algebraic}
\externaldocument[examples-stacks-]{examples-stacks}
\externaldocument[stacks-sheaves-]{stacks-sheaves}
\externaldocument[criteria-]{criteria}
\externaldocument[artin-]{artin}
\externaldocument[quot-]{quot}
\externaldocument[stacks-properties-]{stacks-properties}
\externaldocument[stacks-morphisms-]{stacks-morphisms}
\externaldocument[stacks-limits-]{stacks-limits}
\externaldocument[stacks-cohomology-]{stacks-cohomology}
\externaldocument[stacks-perfect-]{stacks-perfect}
\externaldocument[stacks-introduction-]{stacks-introduction}
\externaldocument[stacks-more-morphisms-]{stacks-more-morphisms}
\externaldocument[stacks-geometry-]{stacks-geometry}
\externaldocument[moduli-]{moduli}
\externaldocument[moduli-curves-]{moduli-curves}
\externaldocument[examples-]{examples}
\externaldocument[exercises-]{exercises}
\externaldocument[guide-]{guide}
\externaldocument[desirables-]{desirables}
\externaldocument[coding-]{coding}
\externaldocument[obsolete-]{obsolete}
\externaldocument[fdl-]{fdl}
\externaldocument[index-]{index}

% Theorem environments.
%
\theoremstyle{plain}
\newtheorem{theorem}[subsection]{Theorem}
\newtheorem{proposition}[subsection]{Proposition}
\newtheorem{lemma}[subsection]{Lemma}

\theoremstyle{definition}
\newtheorem{definition}[subsection]{Definition}
\newtheorem{example}[subsection]{Example}
\newtheorem{exercise}[subsection]{Exercise}
\newtheorem{situation}[subsection]{Situation}

\theoremstyle{remark}
\newtheorem{remark}[subsection]{Remark}
\newtheorem{remarks}[subsection]{Remarks}

\numberwithin{equation}{subsection}

% Macros
%
\def\lim{\mathop{\mathrm{lim}}\nolimits}
\def\colim{\mathop{\mathrm{colim}}\nolimits}
\def\Spec{\mathop{\mathrm{Spec}}}
\def\Hom{\mathop{\mathrm{Hom}}\nolimits}
\def\Ext{\mathop{\mathrm{Ext}}\nolimits}
\def\SheafHom{\mathop{\mathcal{H}\!\mathit{om}}\nolimits}
\def\SheafExt{\mathop{\mathcal{E}\!\mathit{xt}}\nolimits}
\def\Sch{\mathit{Sch}}
\def\Mor{\mathop{\mathrm{Mor}}\nolimits}
\def\Ob{\mathop{\mathrm{Ob}}\nolimits}
\def\Sh{\mathop{\mathit{Sh}}\nolimits}
\def\NL{\mathop{N\!L}\nolimits}
\def\CH{\mathop{\mathrm{CH}}\nolimits}
\def\proetale{{pro\text{-}\acute{e}tale}}
\def\etale{{\acute{e}tale}}
\def\QCoh{\mathit{QCoh}}
\def\Ker{\mathop{\mathrm{Ker}}}
\def\Im{\mathop{\mathrm{Im}}}
\def\Coker{\mathop{\mathrm{Coker}}}
\def\Coim{\mathop{\mathrm{Coim}}}

% Boxtimes
%
\DeclareMathSymbol{\boxtimes}{\mathbin}{AMSa}{"02}

%
% Macros for moduli stacks/spaces
%
\def\QCohstack{\mathcal{QC}\!\mathit{oh}}
\def\Cohstack{\mathcal{C}\!\mathit{oh}}
\def\Spacesstack{\mathcal{S}\!\mathit{paces}}
\def\Quotfunctor{\mathrm{Quot}}
\def\Hilbfunctor{\mathrm{Hilb}}
\def\Curvesstack{\mathcal{C}\!\mathit{urves}}
\def\Polarizedstack{\mathcal{P}\!\mathit{olarized}}
\def\Complexesstack{\mathcal{C}\!\mathit{omplexes}}
% \Pic is the operator that assigns to X its picard group, usage \Pic(X)
% \Picardstack_{X/B} denotes the Picard stack of X over B
% \Picardfunctor_{X/B} denotes the Picard functor of X over B
\def\Pic{\mathop{\mathrm{Pic}}\nolimits}
\def\Picardstack{\mathcal{P}\!\mathit{ic}}
\def\Picardfunctor{\mathrm{Pic}}
\def\Deformationcategory{\mathcal{D}\!\mathit{ef}}
\usepackage{cleveref}
\crefname{originalsection}{original section}{original sections}
\Crefname{originalsection}{Original section}{Original sections}
\makeatletter
\newlabel{derived-example-key-postnikov}{{0D8Z}{source}{}{sourcepointer-1}{}}
\newlabel{derived-example-key-postnikov@cref}{{[originalsection][0][]0D8Z}{[1][0][]source}{}{}{}}
\makeatother
\crefname{originallemma}{original lemma}{original lemmas}
\Crefname{originallemma}{Original lemma}{Original lemmas}
\makeatletter
\newlabel{derived-lemma-existence-postnikov-system}{{0D83}{source}{}{sourcepointer-2}{}}
\newlabel{derived-lemma-existence-postnikov-system@cref}{{[originallemma][0][]0D83}{[1][0][]source}{}{}{}}
\makeatother
\crefname{originalreference}{original reference}{original references}
\Crefname{originalreference}{Original reference}{Original references}
\makeatletter
\newlabel{lemma-abelian-postnikov}{{0D9I}{source}{}{sourcepointer-3}{}}
\newlabel{lemma-abelian-postnikov@cref}{{[originalreference][0][]0D9I}{[1][0][]source}{}{}{}}
\makeatother
\crefname{originalreference}{original reference}{original references}
\Crefname{originalreference}{Original reference}{Original references}
\makeatletter
\newlabel{lemma-simplicial-resolution-Z-site}{{09WI}{source}{}{sourcepointer-4}{}}
\newlabel{lemma-simplicial-resolution-Z-site@cref}{{[originalreference][0][]09WI}{[1][0][]source}{}{}{}}
\makeatother
\crefname{originalreference}{original reference}{original references}
\Crefname{originalreference}{Original reference}{Original references}
\makeatletter
\newlabel{situation-simplicial-site}{{09WE}{source}{}{sourcepointer-5}{}}
\newlabel{situation-simplicial-site@cref}{{[originalreference][0][]09WE}{[1][0][]source}{}{}{}}
\makeatother
\title{Cartesian simplicial systems of the derived category with vanishing self-extensions come from cartesian objects}
\author{The Stacks Project Authors}
\date{stacks-project:spaces-simplicial.tex:local-snapshot-20261012; GNU FDL 1.2 or later}
\begin{document}
\maketitle
\noindent\textit{Editorial note.} All delivered mathematical text is the author's own, copied from the pinned Stacks Project chapter file sources/stacks-spaces-simplicial/pinned-original-source.tex; no statement or proof step was rewritten, reordered or removed.\par\medskip
\noindent\textit{Editorial note.} The counted proof is the author's complete argument: the two 'we omit the computation' spots (source lines 3207: vanishing of the compositions of the constructed complex; line 3257: commutativity of the squares in the comparison diagram) are auxiliary checks sitting mid-proof, not at the main conclusion, which is fully argued; the comparison towers Y'\_\{m,n\} are the key Postnikov example tensored with K\_m; the left adjoints g\_\{n!\} and the component formula used throughout are those established by the restriction-to-components lemma delivered as card L46149 of this batch.\par\medskip
\noindent\textit{Editorial note.} The label lemma-cartesian-object-derived-from-simplicial names the counted statement, and the proof environment immediately following it in the source is delivered verbatim as its proof.\par\medskip
\noindent\textit{Editorial note.} The counted statement is carried by the author's own lemma environment in the chapter source.\par\medskip
\noindent\textit{Allowed standard prerequisites (editorial).}\begin{itemize}
\item The left adjoints g\_n! on the simplicial total site and their component formula. Same chapter, \textbackslash{}label\{lemma-restriction-to-components-site\} (card L46149 of this batch). Lines 3187-3194 and 3253-3261: X\_n = g\_\{n!\}K\_n is formed with the left adjoint of g\_n\textasciicircum{}\{-1\}, and the Hom computations use g\_m\textasciicircum{}\{-1\}g\_\{n!\}K\_n = the direct sum of f\_phi\textasciicircum{}\{-1\}K\_n over phi : [n] -> [m].
\item Derived Categories, existence and comparison of Postnikov systems. Derived Categories chapter, \textbackslash{}label\{derived-lemma-existence-postnikov-system\}. Lines 3224-3229: extends the complex to a Postnikov system when the Hom groups vanish; lines 3296-3301: the second part extends the vertical maps to an isomorphism of Postnikov systems.
\item The simplicial resolution of Z on the site. Same chapter, \textbackslash{}label\{lemma-simplicial-resolution-Z-site\}. Lines 3209-3212: the alternating-sum maps mirror this lemma's construction; lines 3275-3278: the bottom-row arrows of the comparison diagram are its maps.
\item The key Postnikov system computing a homotopy colimit. Derived Categories chapter, \textbackslash{}label\{derived-example-key-postnikov\}; Same chapter, \textbackslash{}label\{lemma-abelian-postnikov\}. Lines 3284-3294: the explicit towers Y'\_\{m,n\} = (direct sums of Z)[-n] tensor K\_m with hocolim K\_m.
\item Uniqueness of homotopy colimits in the derived category. Standard property of homotopy colimits in the derived category. Lines 3302-3310: gamma\_m : g\_m\textasciicircum{}\{-1\}K -> K\_m is assembled from the compatible gamma\_\{m,n\} by uniqueness of hocolims.
\end{itemize}\medskip
\ifdefined\hypertarget\hypertarget{sourcepointer-1}{}\fi
\noindent\textit{Source reference.} Original section 0D8Z: sources/stacks-chow/tags.txt, source lines 16609--16609. Prerequisite pointer invoked at proof line(s) 3268; it names a standard fact of a sibling chapter that the counted argument uses, and every step of the counted proof built on it is retained verbatim in the delivered body.\par\medskip
\ifdefined\hypertarget\hypertarget{sourcepointer-2}{}\fi
\noindent\textit{Source reference.} Original lemma 0D83: sources/stacks-chow/tags.txt, source lines 16578--16578. Prerequisite pointer invoked at proof line(s) 3209; it names a standard fact of a sibling chapter that the counted argument uses, and every step of the counted proof built on it is retained verbatim in the delivered body.\par\medskip
\ifdefined\hypertarget\hypertarget{sourcepointer-3}{}\fi
\noindent\textit{Source reference.} Original reference 0D9I: sources/stacks-spaces-simplicial/pinned-original-source.tex, source lines 3002--3002. Prerequisite pointer invoked at proof line(s) 3267; it names a standard fact of the same chapter that the counted argument uses, and every step of the counted proof built on it is retained verbatim in the delivered body.\par\medskip
\ifdefined\hypertarget\hypertarget{sourcepointer-4}{}\fi
\noindent\textit{Source reference.} Original reference 09WI: sources/stacks-spaces-simplicial/pinned-original-source.tex, source lines 1711--1711. Prerequisite pointer invoked at proof line(s) 3202; it names a standard fact of the same chapter that the counted argument uses, and every step of the counted proof built on it is retained verbatim in the delivered body.\par\medskip
\ifdefined\hypertarget\hypertarget{sourcepointer-5}{}\fi
\noindent\textit{Source reference.} Original reference 09WE: sources/stacks-spaces-simplicial/pinned-original-source.tex, source lines 888--888. Prerequisite pointer invoked at proof line(s) 3176; it names a standard fact of the same chapter that the counted argument uses, and every step of the counted proof built on it is retained verbatim in the delivered body.\par\medskip
\begin{lemma}
\label{lemma-cartesian-object-derived-from-simplicial}
In Situation \ref{situation-simplicial-site}. Let
$(K_n, K_\varphi)$ be a simplicial system of the derived category.
Assume
\begin{enumerate}
\item $(K_n, K_\varphi)$ is cartesian,
\item $\Hom(K_i[t], K_i) = 0$ for $i \geq 0$ and $t > 0$.
\end{enumerate}
Then there exists a cartesian object $K$ of $D(\mathcal{C}_{total})$
whose associated simplicial system is isomorphic to $(K_n, K_\varphi)$.
\end{lemma}

\begin{proof}
Set $X_n = g_{n!}K_n$ in $D(\mathcal{C}_{total})$. For each $n \geq 1$
we have
$$
\Hom(X_n, X_{n - 1}) =
\Hom(K_n, g_n^{-1}g_{n - 1!}K_{n - 1}) =
\bigoplus\nolimits_{\varphi : [n - 1] \to [n]}
\Hom(K_n, f_\varphi^{-1}K_{n - 1})
$$
Thus we get a map $X_n \to X_{n - 1}$ corresponding to the
alternating sum of the maps
$K_\varphi^{-1} : K_n \to f_\varphi^{-1}K_{n - 1}$
where $\varphi$ runs over $\delta^n_0, \ldots, \delta^n_n$.
We can do this because $K_\varphi$ is invertible by assumption (1).
Please observe the similarity with the definition of the maps
in the proof of Lemma \ref{lemma-simplicial-resolution-Z-site}.
We obtain a complex
$$
\ldots \to X_2 \to X_1 \to X_0
$$
in $D(\mathcal{C}_{total})$. We omit the computation which shows
that the compositions are zero. By
Derived Categories, Lemma \ref{derived-lemma-existence-postnikov-system}
if we have
$$
\Hom(X_i[i - j - 2], X_j) = 0\text{ for }i > j + 2
$$
then we can extend this complex to a Postnikov system.
The group is equal to
$$
\Hom(K_i[i - j - 2], g_i^{-1}g_{j!}K_j)
$$
Again using that $(K_n, K_\varphi)$ is cartesian we see that
$g_i^{-1}g_{j!}K_j$ is isomorphic to a finite direct sum of copies of
$K_i$. Hence the group vanishes by assumption (2).
Let the Postnikov system be given by $Y_0 = X_0$ and distinguished
sequences $Y_n \to X_n \to Y_{n - 1} \to Y_n[1]$ for $n \geq 1$.
We set
$$
K = \text{hocolim} Y_n[n]
$$
To finish the proof we have to show that $g_m^{-1}K$ is isomorphic
to $K_m$ for all $m$ compatible with the maps $K_\varphi$. Observe that
$$
g_m^{-1} K = \text{hocolim} g_m^{-1}Y_n[n]
$$
and that $g_m^{-1}Y_n[n]$ is a Postnikov system for $g_m^{-1}X_n$.
Consider the isomorphisms
$$
g_m^{-1}X_n =
\bigoplus\nolimits_{\varphi : [n] \to [m]} f_\varphi^{-1}K_n
\xrightarrow{\bigoplus K_\varphi}
\bigoplus\nolimits_{\varphi : [n] \to [m]} K_m
$$
These maps define an isomorphism of complexes
$$
\xymatrix{
\ldots \ar[r] &
g_m^{-1}X_2 \ar[r] \ar[d] &
g_m^{-1}X_1 \ar[r] \ar[d] &
g_m^{-1}X_0 \ar[d] \\
\ldots \ar[r] &
\bigoplus\nolimits_{\varphi : [2] \to [m]} K_m \ar[r] &
\bigoplus\nolimits_{\varphi : [1] \to [m]} K_m \ar[r] &
\bigoplus\nolimits_{\varphi : [0] \to [m]} K_m
}
$$
in $D(\mathcal{C}_m)$ where the arrows in the bottom row are as
in the proof of Lemma \ref{lemma-simplicial-resolution-Z-site}.
The squares commute by our choice of the arrows of the complex
$\ldots \to X_2 \to X_1 \to X_0$; we omit the computation.
The bottom row complex has a postnikov tower given by
$$
Y'_{m, n} =
\left(\bigoplus\nolimits_{\varphi : [n] \to [m]} \mathbf{Z} \to
\ldots \to
\bigoplus\nolimits_{\varphi : [0] \to [m]} \mathbf{Z}\right)[-n]
\otimes^\mathbf{L}_\mathbf{Z} K_m
$$
and $\text{hocolim} Y'_{m, n} = K_m$
(please compare with the proof of Lemma \ref{lemma-abelian-postnikov}
and Derived Categories, Example \ref{derived-example-key-postnikov}).
Applying the second part of
Derived Categories, Lemma \ref{derived-lemma-existence-postnikov-system}
the vertical maps in the big diagram extend to an isomorphism
of Postnikov systems provided we have
$$
\Hom(g_m^{-1}X_i[i - j - 1], \bigoplus\nolimits_{\varphi : [j] \to [m]} K_m)
= 0\text{ for }i > j + 1
$$
The is true if $\Hom(K_m[i - j - 1], K_m) = 0$ for $i > j + 1$
which holds by assumption (2). Choose an isomorphism given
by $\gamma_{m, n} : g_m^{-1}Y_n \to Y'_{m, n}$ of Postnikov systems
in $D(\mathcal{C}_m)$. By uniqueness of homotopy colimits,
we can find an isomorphism
$$
g_m^{-1} K = \text{hocolim} g_m^{-1}Y_n[n]
\xrightarrow{\gamma_m}
\text{hocolim} Y'_{m, n} = K_m
$$
compatible with $\gamma_{m, n}$.

\medskip\noindent
We still have to prove that the maps $\gamma_m$ fit into commutative diagrams
$$
\xymatrix{
f_\varphi^{-1}g_m^{-1}K \ar[d]_{f_\varphi^{-1}\gamma_m} \ar[r]_{K(\varphi)} &
g_n^{-1}K \ar[d]^{\gamma_n} \\
f_\varphi^{-1}K_m \ar[r]^{K_\varphi} &
K_n
}
$$
for every $\varphi : [m] \to [n]$. Consider the diagram
$$
\xymatrix{
f_\varphi^{-1}(\bigoplus_{\psi : [0] \to [m]} f_\psi^{-1}K_0)
\ar@{=}[r] \ar[d]_{f_\varphi^{-1}(\bigoplus K_\psi)} &
f_\varphi^{-1}g_m^{-1}X_0 \ar[d] \ar[r]_{X_0(\varphi)} &
g_n^{-1}X_0 \ar[d] &
\bigoplus_{\chi : [0] \to [n]} f_\chi^{-1}K_0
\ar@{=}[l] \ar[d]^{\bigoplus K_\chi} \\
f_\varphi^{-1}(\bigoplus_{\psi : [0] \to [m]} K_m) \ar@{=}[d] &
f_\varphi^{-1}g_m^{-1}K \ar[d]_{f_\varphi^{-1}\gamma_m} \ar[r]_{K(\varphi)} &
g_n^{-1}K \ar[d]^{\gamma_n} &
\bigoplus_{\chi : [0] \to [n]} K_n \ar@{=}[d] \\
f_\varphi^{-1}Y'_{0, m} \ar[r] &
f_\varphi^{-1}K_m \ar[r]^{K_\varphi} &
K_n &
Y'_{0, n} \ar[l]
}
$$
The top middle square is commutative as $X_0 \to K$ is a morphism
of simplicial objects. The left, resp.\ the right rectangles are
commutative as $\gamma_m$, resp.\ $\gamma_n$ is compatible with
$\gamma_{0, m}$, resp.\ $\gamma_{0, n}$ which are the arrows
$\bigoplus K_\psi$ and $\bigoplus K_\chi$ in the diagram.
Going around the outer rectangle of the diagram
is commutative as $(K_n, K_\varphi)$ is a simplical system
and the map $X_0(\varphi)$ is given by the obvious identifications
$f_\varphi^{-1}f_\psi^{-1}K_0 = f_{\varphi \circ \psi}^{-1}K_0$.
Note that the arrow $\bigoplus_\psi K_m \to Y'_{0, m} \to K_m$
induces an isomorphism on any of the direct summands
(because of our explicit construction of the Postnikov
systems $Y'_{i, j}$ above).
Hence, if we take a direct summand of
the upper left and corner, then this maps isomorphically to
$f_\varphi^{-1}g_m^{-1}K$ as $\gamma_m$ is an isomorphism.
Working out what the above says,
but looking only at this direct summand we conclude the lower
middle square commutes as we well. This concludes the proof.
\end{proof}

\end{document}
